
\documentclass{article}
\usepackage{colortbl}
\usepackage{makecell}
\usepackage{multirow}
\usepackage{supertabular}

\begin{document}

\newcounter{utterance}

\centering \large Interaction Transcript for game `clean\_up', experiment `1\_medium\_7obj\_uk', episode 0 with gpt{-}5.4{-}azure.
\vspace{24pt}

{ \footnotesize  \setcounter{utterance}{1}
\setlength{\tabcolsep}{0pt}
\begin{supertabular}{c@{$\;$}|p{.15\linewidth}@{}p{.15\linewidth}p{.15\linewidth}p{.15\linewidth}p{.15\linewidth}p{.15\linewidth}}
    \# & $\;$A & \multicolumn{4}{c}{Game Master} & $\;\:$B\\
    \hline

    \theutterance \stepcounter{utterance}  
    & & \multicolumn{4}{p{0.6\linewidth}}{
        \cellcolor[rgb]{0.9,0.9,0.9}{
            \makecell[{{p{\linewidth}}}]{
                \texttt{\tiny{[P1$\langle$GM]}}
                \texttt{Я ваш ведучий гри, і ви граєте у спільну гру з наступною сіткою як ігровим полем:} \\
\\ 
\texttt{```} \\
\texttt{ 1234567} \\
\texttt{╔════╤══╗} \\
\texttt{║◌◌◌O│PM║ 1} \\
\texttt{║◌◌T◌├──╢ 2} \\
\texttt{║◌◌◌◌│◌◌║ 3} \\
\texttt{║S┌──┼──╢ 4} \\
\texttt{║◌│◌A│◌◌║ 5} \\
\texttt{╟─┴┬─┤◌◌║ 6} \\
\texttt{║◌◌│◌│D◌║ 7} \\
\texttt{╚══╧═╧══╝} \\
\\ 
\texttt{```} \\
\\ 
\texttt{* Верхній край показує координати x, що зростають праворуч, а правий край — координати y, що зростають донизу.} \\
\texttt{* Наступні об'єкти випадково розміщені на вашій сітці: 'P', 'O', 'T', 'S', 'D', 'A', 'M'.} \\
\\ 
\\ 
\texttt{Інший гравець бачить варіант ігрового поля, де об'єкти розміщені в інших випадкових місцях. Ви не бачите поле іншого гравця, а він не бачить вашого.} \\
\\ 
\\ 
\texttt{**Мета:**} \\
\\ 
\texttt{Обидва гравці повинні переміщати об'єкти на своєму фоні так, щоб однакові об'єкти опинилися в однакових координатах. Ви повинні спілкуватися з іншим гравцем, щоб узгодити спільну цільову конфігурацію.} \\
\\ 
\\ 
\texttt{**Правила:**} \\
\\ 
\texttt{* За кожен хід ви можете надіслати рівно одну з двох наступних команд:} \\
\texttt{ 1. `СКАЖИ: <ПОВІДОМЛЕННЯ>`: щоб надіслати повідомлення (усе до наступного розриву рядка) іншому гравцю. Я передам його вашому партнеру.} \\
\texttt{2. `ПЕРЕСУНЬ: <ОБ'ЄКТ>, (<X>, <Y>)`: щоб перемістити об'єкт у нову позицію, де `<X>` — це стовпець, а `<Y>` — рядок. Я повідомлю вас, чи був ваш хід успішним.} \\
\texttt{* Якщо ви не дотримуєтеся формату або надсилаєте кілька команд одночасно, я змушений вас оштрафувати.} \\
\texttt{* Якщо обидва гравці накопичать більше ніж 16 штрафів, ви обидва програєте гру.} \\
\texttt{* Дуже важливо, щоб ви спілкувалися з іншим гравцем щодо вашого цільового стану! *Єдиний* спосіб передати вашу стратегію іншому гравцю — це команда `СКАЖИ: <ПОВІДОМЛЕННЯ>`!} \\
\\ 
\\ 
\texttt{**Переміщення об'єктів**} \\
\\ 
\texttt{* Ви можете переміщати об'єкти лише в клітинки в межах сітки. Цільова клітинка повинна бути порожньою, тобто містити лише символ '◌'.} \\
\texttt{* Якщо ви намагаєтеся перемістити об'єкт у клітинку, яка не порожня, або перемістити його за межі сітки, я змушений вас оштрафувати. Ви отримаєте ще одну спробу.} \\
\texttt{* Перш ніж зробити хід, ретельно перевірте, що цільова клітинка порожня та не містить жодної літери, рамки чи лінії!} \\
\\ 
\\ 
\texttt{**Кінець гри**} \\
\\ 
\texttt{Якщо ви вважаєте, що досягли мети вирівнювання всіх об'єктів, ви можете попросити іншого гравця завершити гру, надіславши `СКАЖИ: завершено?`. Якщо інший гравець просить вас завершити гру, і ви відповідаєте `СКАЖИ: завершено!`, гра завершиться.} \\
\\ 
\texttt{Обидва гравці перемагають, якщо гра завершується протягом 28 раундів, де один раунд визначається як надсилання обома гравцями по одній дійсній команді.} \\
\\ 
\texttt{**Підрахунок очок:**} \\
\\ 
\texttt{Чим ближче однакові об'єкти на обох ігрових полях, тим більше очок ви отримуєте. Штрафи зменшують ваші очки. Чи зможете ви побити рекорд?} \\
\\ 
\\ 
\texttt{Будь ласка, надішліть повідомлення іншому гравцю, щоб почати гру!} \\
            }
        }
    }
    & \\ \\

    \theutterance \stepcounter{utterance}  
    & \multicolumn{4}{p{0.6\linewidth}}{
        \cellcolor[rgb]{0.9,0.9,0.9}{
            \makecell[{{p{\linewidth}}}]{
                \texttt{\tiny{[P1$\rangle$GM]}}
                \texttt{СКАЖИ: Пропоную ціль = моя поточна конфігурація: O(4,1), P(6,1), M(7,1), T(3,2), S(1,4), A(4,5), D(6,7). Якщо згоден, рухай свої об’єкти до цих координат.} \\
            }
        }
    }
    & & \\ \\

    \theutterance \stepcounter{utterance}  
    & & & \multicolumn{2}{p{0.3\linewidth}}{
        \cellcolor[rgb]{0.9,0.9,0.9}{
            \makecell[{{p{\linewidth}}}]{
                \texttt{\tiny{[GM$|$GM]}}
                \texttt{СКАЖИ: Пропоную ціль = моя поточна конфігурація: O(4,1), P(6,1), M(7,1), T(3,2), S(1,4), A(4,5), D(6,7). Якщо згоден, рухай свої об’єкти до цих координат.} \\
            }
        }
    }
    & & \\ \\

    \theutterance \stepcounter{utterance}  
    & & \multicolumn{4}{p{0.6\linewidth}}{
        \cellcolor[rgb]{0.9,0.9,0.9}{
            \makecell[{{p{\linewidth}}}]{
                \texttt{\tiny{[GM$\rangle$P2]}}
                \texttt{Я ваш ведучий гри, і ви граєте у спільну гру з наступною сіткою як ігровим полем:} \\
\\ 
\texttt{```} \\
\texttt{ 1234567} \\
\texttt{╔════╤══╗} \\
\texttt{║◌◌◌◌│O◌║ 1} \\
\texttt{║P◌M◌├──╢ 2} \\
\texttt{║◌S◌◌│T◌║ 3} \\
\texttt{║◌┌──┼──╢ 4} \\
\texttt{║A│◌◌│◌◌║ 5} \\
\texttt{╟─┴┬─┤◌D║ 6} \\
\texttt{║◌◌│◌│◌◌║ 7} \\
\texttt{╚══╧═╧══╝} \\
\\ 
\texttt{```} \\
\\ 
\texttt{* Верхній край показує координати x, що зростають праворуч, а правий край — координати y, що зростають донизу.} \\
\texttt{* Наступні об'єкти випадково розміщені на вашій сітці: 'P', 'O', 'T', 'S', 'D', 'A', 'M'.} \\
\\ 
\\ 
\texttt{Інший гравець бачить варіант ігрового поля, де об'єкти розміщені в інших випадкових місцях. Ви не бачите поле іншого гравця, а він не бачить вашого.} \\
\\ 
\\ 
\texttt{**Мета:**} \\
\\ 
\texttt{Обидва гравці повинні переміщати об'єкти на своєму фоні так, щоб однакові об'єкти опинилися в однакових координатах. Ви повинні спілкуватися з іншим гравцем, щоб узгодити спільну цільову конфігурацію.} \\
\\ 
\\ 
\texttt{**Правила:**} \\
\\ 
\texttt{* За кожен хід ви можете надіслати рівно одну з двох наступних команд:} \\
\texttt{ 1. `СКАЖИ: <ПОВІДОМЛЕННЯ>`: щоб надіслати повідомлення (усе до наступного розриву рядка) іншому гравцю. Я передам його вашому партнеру.} \\
\texttt{2. `ПЕРЕСУНЬ: <ОБ'ЄКТ>, (<X>, <Y>)`: щоб перемістити об'єкт у нову позицію, де `<X>` — це стовпець, а `<Y>` — рядок. Я повідомлю вас, чи був ваш хід успішним.} \\
\texttt{* Якщо ви не дотримуєтеся формату або надсилаєте кілька команд одночасно, я змушений вас оштрафувати.} \\
\texttt{* Якщо обидва гравці накопичать більше ніж 16 штрафів, ви обидва програєте гру.} \\
\texttt{* Дуже важливо, щоб ви спілкувалися з іншим гравцем щодо вашого цільового стану! *Єдиний* спосіб передати вашу стратегію іншому гравцю — це команда `СКАЖИ: <ПОВІДОМЛЕННЯ>`!} \\
\\ 
\\ 
\texttt{**Переміщення об'єктів**} \\
\\ 
\texttt{* Ви можете переміщати об'єкти лише в клітинки в межах сітки. Цільова клітинка повинна бути порожньою, тобто містити лише символ '◌'.} \\
\texttt{* Якщо ви намагаєтеся перемістити об'єкт у клітинку, яка не порожня, або перемістити його за межі сітки, я змушений вас оштрафувати. Ви отримаєте ще одну спробу.} \\
\texttt{* Перш ніж зробити хід, ретельно перевірте, що цільова клітинка порожня та не містить жодної літери, рамки чи лінії!} \\
\\ 
\\ 
\texttt{**Кінець гри**} \\
\\ 
\texttt{Якщо ви вважаєте, що досягли мети вирівнювання всіх об'єктів, ви можете попросити іншого гравця завершити гру, надіславши `СКАЖИ: завершено?`. Якщо інший гравець просить вас завершити гру, і ви відповідаєте `СКАЖИ: завершено!`, гра завершиться.} \\
\\ 
\texttt{Обидва гравці перемагають, якщо гра завершується протягом 28 раундів, де один раунд визначається як надсилання обома гравцями по одній дійсній команді.} \\
\\ 
\texttt{**Підрахунок очок:**} \\
\\ 
\texttt{Чим ближче однакові об'єкти на обох ігрових полях, тим більше очок ви отримуєте. Штрафи зменшують ваші очки. Чи зможете ви побити рекорд?} \\
\\ 
\texttt{Інший гравець почав гру, надіславши це повідомлення:} \\
\\ 
\texttt{"Пропоную ціль = моя поточна конфігурація: O(4,1), P(6,1), M(7,1), T(3,2), S(1,4), A(4,5), D(6,7). Якщо згоден, рухай свої об’єкти до цих координат."} \\
\\ 
\texttt{Яка ваша перша команда?} \\
            }
        }
    }
    & \\ \\

    \theutterance \stepcounter{utterance}  
    & & & \multicolumn{4}{p{0.6\linewidth}}{
        \cellcolor[rgb]{0.9,0.9,0.9}{
            \makecell[{{p{\linewidth}}}]{
                \texttt{\tiny{[GM$\langle$P2]}}
                \texttt{СКАЖИ: Згоден, приймаю твою конфігурацію як ціль: O(4,1), P(6,1), M(7,1), T(3,2), S(1,4), A(4,5), D(6,7). Рухаю свої об’єкти туди.} \\
            }
        }
    }
     \\ \\

    \theutterance \stepcounter{utterance}  
    & & & \multicolumn{2}{p{0.3\linewidth}}{
        \cellcolor[rgb]{0.9,0.9,0.9}{
            \makecell[{{p{\linewidth}}}]{
                \texttt{\tiny{[GM$|$GM]}}
                \texttt{СКАЖИ: Згоден, приймаю твою конфігурацію як ціль: O(4,1), P(6,1), M(7,1), T(3,2), S(1,4), A(4,5), D(6,7). Рухаю свої об’єкти туди.} \\
            }
        }
    }
    & & \\ \\

    \theutterance \stepcounter{utterance}  
    & & \multicolumn{4}{p{0.6\linewidth}}{
        \cellcolor[rgb]{0.9,0.9,0.9}{
            \makecell[{{p{\linewidth}}}]{
                \texttt{\tiny{[P1$\langle$GM]}}
                \texttt{Ваше повідомлення передано іншому гравцю.} \\
\\ 
\\ 
\texttt{Зараз ви граєте раунд 1 з максимум 28.} \\
\texttt{Разом ви накопичили 0 з 16 штрафів.} \\
\texttt{Інший гравець надіслав це повідомлення:} \\
\\ 
\\ 
\texttt{"Згоден, приймаю твою конфігурацію як ціль: O(4,1), P(6,1), M(7,1), T(3,2), S(1,4), A(4,5), D(6,7). Рухаю свої об’єкти туди."} \\
\\ 
\\ 
\texttt{Яка ваша наступна команда?} \\
            }
        }
    }
    & \\ \\

    \theutterance \stepcounter{utterance}  
    & \multicolumn{4}{p{0.6\linewidth}}{
        \cellcolor[rgb]{0.9,0.9,0.9}{
            \makecell[{{p{\linewidth}}}]{
                \texttt{\tiny{[P1$\rangle$GM]}}
                \texttt{СКАЖИ: Моє поле вже відповідає цілі, тож я не рухаюсь. Коли виставиш усі об’єкти в ці координати, надішли "завершено?".} \\
            }
        }
    }
    & & \\ \\

    \theutterance \stepcounter{utterance}  
    & & & \multicolumn{2}{p{0.3\linewidth}}{
        \cellcolor[rgb]{0.9,0.9,0.9}{
            \makecell[{{p{\linewidth}}}]{
                \texttt{\tiny{[GM$|$GM]}}
                \texttt{СКАЖИ: Моє поле вже відповідає цілі, тож я не рухаюсь. Коли виставиш усі об’єкти в ці координати, надішли "завершено?".} \\
            }
        }
    }
    & & \\ \\

    \theutterance \stepcounter{utterance}  
    & & \multicolumn{4}{p{0.6\linewidth}}{
        \cellcolor[rgb]{0.9,0.9,0.9}{
            \makecell[{{p{\linewidth}}}]{
                \texttt{\tiny{[GM$\rangle$P2]}}
                \texttt{Ваше повідомлення передано іншому гравцю.} \\
\\ 
\\ 
\texttt{Зараз ви граєте раунд 2 з максимум 28.} \\
\texttt{Разом ви накопичили 0 з 16 штрафів.} \\
\texttt{Інший гравець надіслав це повідомлення:} \\
\\ 
\\ 
\texttt{"Моє поле вже відповідає цілі, тож я не рухаюсь. Коли виставиш усі об’єкти в ці координати, надішли "завершено?"."} \\
\\ 
\\ 
\texttt{Яка ваша наступна команда?} \\
            }
        }
    }
    & \\ \\

    \theutterance \stepcounter{utterance}  
    & & & \multicolumn{4}{p{0.6\linewidth}}{
        \cellcolor[rgb]{0.9,0.9,0.9}{
            \makecell[{{p{\linewidth}}}]{
                \texttt{\tiny{[GM$\langle$P2]}}
                \texttt{ПЕРЕСУНЬ: M, (7, 1)} \\
            }
        }
    }
     \\ \\

    \theutterance \stepcounter{utterance}  
    & & & \multicolumn{2}{p{0.3\linewidth}}{
        \cellcolor[rgb]{0.9,0.9,0.9}{
            \makecell[{{p{\linewidth}}}]{
                \texttt{\tiny{[GM$|$GM]}}
                \texttt{ПЕРЕСУНЬ: M, (7, 1)} \\
            }
        }
    }
    & & \\ \\

    \theutterance \stepcounter{utterance}  
    & & & \multicolumn{2}{p{0.3\linewidth}}{
        \cellcolor[rgb]{0.9,0.9,0.9}{
            \makecell[{{p{\linewidth}}}]{
                \texttt{\tiny{[GM$|$GM]}}
                \texttt{M moved to (7, 1)} \\
            }
        }
    }
    & & \\ \\

    \theutterance \stepcounter{utterance}  
    & & \multicolumn{4}{p{0.6\linewidth}}{
        \cellcolor[rgb]{0.9,0.9,0.9}{
            \makecell[{{p{\linewidth}}}]{
                \texttt{\tiny{[P1$\langle$GM]}}
                \texttt{Ваше повідомлення передано іншому гравцю.} \\
\\ 
\\ 
\texttt{Зараз ви граєте раунд 2 з максимум 28.} \\
\texttt{Разом ви накопичили 0 з 16 штрафів.} \\
\texttt{Інший гравець перемістив об'єкт на своїй сітці.} \\
\\ 
\\ 
\texttt{Яка ваша наступна команда?} \\
            }
        }
    }
    & \\ \\

    \theutterance \stepcounter{utterance}  
    & \multicolumn{4}{p{0.6\linewidth}}{
        \cellcolor[rgb]{0.9,0.9,0.9}{
            \makecell[{{p{\linewidth}}}]{
                \texttt{\tiny{[P1$\rangle$GM]}}
                \texttt{СКАЖИ: Продовжуй до цілі. Якщо хочеш, можеш надсилати короткі оновлення по кроках; коли все буде готово, напиши "завершено?".} \\
            }
        }
    }
    & & \\ \\

    \theutterance \stepcounter{utterance}  
    & & & \multicolumn{2}{p{0.3\linewidth}}{
        \cellcolor[rgb]{0.9,0.9,0.9}{
            \makecell[{{p{\linewidth}}}]{
                \texttt{\tiny{[GM$|$GM]}}
                \texttt{СКАЖИ: Продовжуй до цілі. Якщо хочеш, можеш надсилати короткі оновлення по кроках; коли все буде готово, напиши "завершено?".} \\
            }
        }
    }
    & & \\ \\

    \theutterance \stepcounter{utterance}  
    & & \multicolumn{4}{p{0.6\linewidth}}{
        \cellcolor[rgb]{0.9,0.9,0.9}{
            \makecell[{{p{\linewidth}}}]{
                \texttt{\tiny{[GM$\rangle$P2]}}
                \texttt{Об'єкт 'M' успішно переміщено в (7,1). Ваша оновлена сітка виглядає так:} \\
\texttt{```} \\
\texttt{ 1234567} \\
\texttt{╔════╤══╗} \\
\texttt{║◌◌◌◌│OM║ 1} \\
\texttt{║P◌◌◌├──╢ 2} \\
\texttt{║◌S◌◌│T◌║ 3} \\
\texttt{║◌┌──┼──╢ 4} \\
\texttt{║A│◌◌│◌◌║ 5} \\
\texttt{╟─┴┬─┤◌D║ 6} \\
\texttt{║◌◌│◌│◌◌║ 7} \\
\texttt{╚══╧═╧══╝} \\
\texttt{```} \\
\texttt{Зараз ви граєте раунд 3 з максимум 28.} \\
\texttt{Разом ви накопичили 0 з 16 штрафів.} \\
\texttt{Інший гравець надіслав це повідомлення:} \\
\\ 
\\ 
\texttt{"Продовжуй до цілі. Якщо хочеш, можеш надсилати короткі оновлення по кроках; коли все буде готово, напиши "завершено?"."} \\
\\ 
\\ 
\texttt{Яка ваша наступна команда?} \\
            }
        }
    }
    & \\ \\

    \theutterance \stepcounter{utterance}  
    & & & \multicolumn{4}{p{0.6\linewidth}}{
        \cellcolor[rgb]{0.9,0.9,0.9}{
            \makecell[{{p{\linewidth}}}]{
                \texttt{\tiny{[GM$\langle$P2]}}
                \texttt{ПЕРЕСУНЬ: O, (4, 1)} \\
            }
        }
    }
     \\ \\

    \theutterance \stepcounter{utterance}  
    & & & \multicolumn{2}{p{0.3\linewidth}}{
        \cellcolor[rgb]{0.9,0.9,0.9}{
            \makecell[{{p{\linewidth}}}]{
                \texttt{\tiny{[GM$|$GM]}}
                \texttt{ПЕРЕСУНЬ: O, (4, 1)} \\
            }
        }
    }
    & & \\ \\

    \theutterance \stepcounter{utterance}  
    & & & \multicolumn{2}{p{0.3\linewidth}}{
        \cellcolor[rgb]{0.9,0.9,0.9}{
            \makecell[{{p{\linewidth}}}]{
                \texttt{\tiny{[GM$|$GM]}}
                \texttt{O moved to (4, 1)} \\
            }
        }
    }
    & & \\ \\

    \theutterance \stepcounter{utterance}  
    & & \multicolumn{4}{p{0.6\linewidth}}{
        \cellcolor[rgb]{0.9,0.9,0.9}{
            \makecell[{{p{\linewidth}}}]{
                \texttt{\tiny{[P1$\langle$GM]}}
                \texttt{Ваше повідомлення передано іншому гравцю.} \\
\\ 
\\ 
\texttt{Зараз ви граєте раунд 3 з максимум 28.} \\
\texttt{Разом ви накопичили 0 з 16 штрафів.} \\
\texttt{Інший гравець перемістив об'єкт на своїй сітці.} \\
\\ 
\\ 
\texttt{Яка ваша наступна команда?} \\
            }
        }
    }
    & \\ \\

    \theutterance \stepcounter{utterance}  
    & \multicolumn{4}{p{0.6\linewidth}}{
        \cellcolor[rgb]{0.9,0.9,0.9}{
            \makecell[{{p{\linewidth}}}]{
                \texttt{\tiny{[P1$\rangle$GM]}}
                \texttt{СКАЖИ: Нагадую цільові координати: O(4,1), P(6,1), M(7,1), T(3,2), S(1,4), A(4,5), D(6,7). Коли все виставиш, напиши "завершено?".} \\
            }
        }
    }
    & & \\ \\

    \theutterance \stepcounter{utterance}  
    & & & \multicolumn{2}{p{0.3\linewidth}}{
        \cellcolor[rgb]{0.9,0.9,0.9}{
            \makecell[{{p{\linewidth}}}]{
                \texttt{\tiny{[GM$|$GM]}}
                \texttt{СКАЖИ: Нагадую цільові координати: O(4,1), P(6,1), M(7,1), T(3,2), S(1,4), A(4,5), D(6,7). Коли все виставиш, напиши "завершено?".} \\
            }
        }
    }
    & & \\ \\

    \theutterance \stepcounter{utterance}  
    & & \multicolumn{4}{p{0.6\linewidth}}{
        \cellcolor[rgb]{0.9,0.9,0.9}{
            \makecell[{{p{\linewidth}}}]{
                \texttt{\tiny{[GM$\rangle$P2]}}
                \texttt{Об'єкт 'O' успішно переміщено в (4,1). Ваша оновлена сітка виглядає так:} \\
\texttt{```} \\
\texttt{ 1234567} \\
\texttt{╔════╤══╗} \\
\texttt{║◌◌◌O│◌M║ 1} \\
\texttt{║P◌◌◌├──╢ 2} \\
\texttt{║◌S◌◌│T◌║ 3} \\
\texttt{║◌┌──┼──╢ 4} \\
\texttt{║A│◌◌│◌◌║ 5} \\
\texttt{╟─┴┬─┤◌D║ 6} \\
\texttt{║◌◌│◌│◌◌║ 7} \\
\texttt{╚══╧═╧══╝} \\
\texttt{```} \\
\texttt{Зараз ви граєте раунд 4 з максимум 28.} \\
\texttt{Разом ви накопичили 0 з 16 штрафів.} \\
\texttt{Інший гравець надіслав це повідомлення:} \\
\\ 
\\ 
\texttt{"Нагадую цільові координати: O(4,1), P(6,1), M(7,1), T(3,2), S(1,4), A(4,5), D(6,7). Коли все виставиш, напиши "завершено?"."} \\
\\ 
\\ 
\texttt{Яка ваша наступна команда?} \\
            }
        }
    }
    & \\ \\

    \theutterance \stepcounter{utterance}  
    & & & \multicolumn{4}{p{0.6\linewidth}}{
        \cellcolor[rgb]{0.9,0.9,0.9}{
            \makecell[{{p{\linewidth}}}]{
                \texttt{\tiny{[GM$\langle$P2]}}
                \texttt{ПЕРЕСУНЬ: P, (6, 1)} \\
            }
        }
    }
     \\ \\

    \theutterance \stepcounter{utterance}  
    & & & \multicolumn{2}{p{0.3\linewidth}}{
        \cellcolor[rgb]{0.9,0.9,0.9}{
            \makecell[{{p{\linewidth}}}]{
                \texttt{\tiny{[GM$|$GM]}}
                \texttt{ПЕРЕСУНЬ: P, (6, 1)} \\
            }
        }
    }
    & & \\ \\

    \theutterance \stepcounter{utterance}  
    & & & \multicolumn{2}{p{0.3\linewidth}}{
        \cellcolor[rgb]{0.9,0.9,0.9}{
            \makecell[{{p{\linewidth}}}]{
                \texttt{\tiny{[GM$|$GM]}}
                \texttt{P moved to (6, 1)} \\
            }
        }
    }
    & & \\ \\

    \theutterance \stepcounter{utterance}  
    & & \multicolumn{4}{p{0.6\linewidth}}{
        \cellcolor[rgb]{0.9,0.9,0.9}{
            \makecell[{{p{\linewidth}}}]{
                \texttt{\tiny{[P1$\langle$GM]}}
                \texttt{Ваше повідомлення передано іншому гравцю.} \\
\\ 
\\ 
\texttt{Зараз ви граєте раунд 4 з максимум 28.} \\
\texttt{Разом ви накопичили 0 з 16 штрафів.} \\
\texttt{Інший гравець перемістив об'єкт на своїй сітці.} \\
\\ 
\\ 
\texttt{Яка ваша наступна команда?} \\
            }
        }
    }
    & \\ \\

    \theutterance \stepcounter{utterance}  
    & \multicolumn{4}{p{0.6\linewidth}}{
        \cellcolor[rgb]{0.9,0.9,0.9}{
            \makecell[{{p{\linewidth}}}]{
                \texttt{\tiny{[P1$\rangle$GM]}}
                \texttt{СКАЖИ: Скільки об’єктів ще лишилось переставити? Якщо вже всі на місці, напиши "завершено?".} \\
            }
        }
    }
    & & \\ \\

    \theutterance \stepcounter{utterance}  
    & & & \multicolumn{2}{p{0.3\linewidth}}{
        \cellcolor[rgb]{0.9,0.9,0.9}{
            \makecell[{{p{\linewidth}}}]{
                \texttt{\tiny{[GM$|$GM]}}
                \texttt{СКАЖИ: Скільки об’єктів ще лишилось переставити? Якщо вже всі на місці, напиши "завершено?".} \\
            }
        }
    }
    & & \\ \\

    \theutterance \stepcounter{utterance}  
    & & \multicolumn{4}{p{0.6\linewidth}}{
        \cellcolor[rgb]{0.9,0.9,0.9}{
            \makecell[{{p{\linewidth}}}]{
                \texttt{\tiny{[GM$\rangle$P2]}}
                \texttt{Об'єкт 'P' успішно переміщено в (6,1). Ваша оновлена сітка виглядає так:} \\
\texttt{```} \\
\texttt{ 1234567} \\
\texttt{╔════╤══╗} \\
\texttt{║◌◌◌O│PM║ 1} \\
\texttt{║◌◌◌◌├──╢ 2} \\
\texttt{║◌S◌◌│T◌║ 3} \\
\texttt{║◌┌──┼──╢ 4} \\
\texttt{║A│◌◌│◌◌║ 5} \\
\texttt{╟─┴┬─┤◌D║ 6} \\
\texttt{║◌◌│◌│◌◌║ 7} \\
\texttt{╚══╧═╧══╝} \\
\texttt{```} \\
\texttt{Зараз ви граєте раунд 5 з максимум 28.} \\
\texttt{Разом ви накопичили 0 з 16 штрафів.} \\
\texttt{Інший гравець надіслав це повідомлення:} \\
\\ 
\\ 
\texttt{"Скільки об’єктів ще лишилось переставити? Якщо вже всі на місці, напиши "завершено?"."} \\
\\ 
\\ 
\texttt{Яка ваша наступна команда?} \\
            }
        }
    }
    & \\ \\

    \theutterance \stepcounter{utterance}  
    & & & \multicolumn{4}{p{0.6\linewidth}}{
        \cellcolor[rgb]{0.9,0.9,0.9}{
            \makecell[{{p{\linewidth}}}]{
                \texttt{\tiny{[GM$\langle$P2]}}
                \texttt{ПЕРЕСУНЬ: T, (3, 2)} \\
            }
        }
    }
     \\ \\

    \theutterance \stepcounter{utterance}  
    & & & \multicolumn{2}{p{0.3\linewidth}}{
        \cellcolor[rgb]{0.9,0.9,0.9}{
            \makecell[{{p{\linewidth}}}]{
                \texttt{\tiny{[GM$|$GM]}}
                \texttt{ПЕРЕСУНЬ: T, (3, 2)} \\
            }
        }
    }
    & & \\ \\

    \theutterance \stepcounter{utterance}  
    & & & \multicolumn{2}{p{0.3\linewidth}}{
        \cellcolor[rgb]{0.9,0.9,0.9}{
            \makecell[{{p{\linewidth}}}]{
                \texttt{\tiny{[GM$|$GM]}}
                \texttt{T moved to (3, 2)} \\
            }
        }
    }
    & & \\ \\

    \theutterance \stepcounter{utterance}  
    & & \multicolumn{4}{p{0.6\linewidth}}{
        \cellcolor[rgb]{0.9,0.9,0.9}{
            \makecell[{{p{\linewidth}}}]{
                \texttt{\tiny{[P1$\langle$GM]}}
                \texttt{Ваше повідомлення передано іншому гравцю.} \\
\\ 
\\ 
\texttt{Зараз ви граєте раунд 5 з максимум 28.} \\
\texttt{Разом ви накопичили 0 з 16 штрафів.} \\
\texttt{Інший гравець перемістив об'єкт на своїй сітці.} \\
\\ 
\\ 
\texttt{Яка ваша наступна команда?} \\
            }
        }
    }
    & \\ \\

    \theutterance \stepcounter{utterance}  
    & \multicolumn{4}{p{0.6\linewidth}}{
        \cellcolor[rgb]{0.9,0.9,0.9}{
            \makecell[{{p{\linewidth}}}]{
                \texttt{\tiny{[P1$\rangle$GM]}}
                \texttt{СКАЖИ: Якщо це був останній або передостанній хід, перевір усі координати ще раз; коли все збігатиметься, напиши "завершено?".} \\
            }
        }
    }
    & & \\ \\

    \theutterance \stepcounter{utterance}  
    & & & \multicolumn{2}{p{0.3\linewidth}}{
        \cellcolor[rgb]{0.9,0.9,0.9}{
            \makecell[{{p{\linewidth}}}]{
                \texttt{\tiny{[GM$|$GM]}}
                \texttt{СКАЖИ: Якщо це був останній або передостанній хід, перевір усі координати ще раз; коли все збігатиметься, напиши "завершено?".} \\
            }
        }
    }
    & & \\ \\

    \theutterance \stepcounter{utterance}  
    & & \multicolumn{4}{p{0.6\linewidth}}{
        \cellcolor[rgb]{0.9,0.9,0.9}{
            \makecell[{{p{\linewidth}}}]{
                \texttt{\tiny{[GM$\rangle$P2]}}
                \texttt{Об'єкт 'T' успішно переміщено в (3,2). Ваша оновлена сітка виглядає так:} \\
\texttt{```} \\
\texttt{ 1234567} \\
\texttt{╔════╤══╗} \\
\texttt{║◌◌◌O│PM║ 1} \\
\texttt{║◌◌T◌├──╢ 2} \\
\texttt{║◌S◌◌│◌◌║ 3} \\
\texttt{║◌┌──┼──╢ 4} \\
\texttt{║A│◌◌│◌◌║ 5} \\
\texttt{╟─┴┬─┤◌D║ 6} \\
\texttt{║◌◌│◌│◌◌║ 7} \\
\texttt{╚══╧═╧══╝} \\
\texttt{```} \\
\texttt{Зараз ви граєте раунд 6 з максимум 28.} \\
\texttt{Разом ви накопичили 0 з 16 штрафів.} \\
\texttt{Інший гравець надіслав це повідомлення:} \\
\\ 
\\ 
\texttt{"Якщо це був останній або передостанній хід, перевір усі координати ще раз; коли все збігатиметься, напиши "завершено?"."} \\
\\ 
\\ 
\texttt{Яка ваша наступна команда?} \\
            }
        }
    }
    & \\ \\

    \theutterance \stepcounter{utterance}  
    & & & \multicolumn{4}{p{0.6\linewidth}}{
        \cellcolor[rgb]{0.9,0.9,0.9}{
            \makecell[{{p{\linewidth}}}]{
                \texttt{\tiny{[GM$\langle$P2]}}
                \texttt{ПЕРЕСУНЬ: S, (1, 4)} \\
            }
        }
    }
     \\ \\

    \theutterance \stepcounter{utterance}  
    & & & \multicolumn{2}{p{0.3\linewidth}}{
        \cellcolor[rgb]{0.9,0.9,0.9}{
            \makecell[{{p{\linewidth}}}]{
                \texttt{\tiny{[GM$|$GM]}}
                \texttt{ПЕРЕСУНЬ: S, (1, 4)} \\
            }
        }
    }
    & & \\ \\

    \theutterance \stepcounter{utterance}  
    & & & \multicolumn{2}{p{0.3\linewidth}}{
        \cellcolor[rgb]{0.9,0.9,0.9}{
            \makecell[{{p{\linewidth}}}]{
                \texttt{\tiny{[GM$|$GM]}}
                \texttt{S moved to (1, 4)} \\
            }
        }
    }
    & & \\ \\

    \theutterance \stepcounter{utterance}  
    & & \multicolumn{4}{p{0.6\linewidth}}{
        \cellcolor[rgb]{0.9,0.9,0.9}{
            \makecell[{{p{\linewidth}}}]{
                \texttt{\tiny{[P1$\langle$GM]}}
                \texttt{Ваше повідомлення передано іншому гравцю.} \\
\\ 
\\ 
\texttt{Зараз ви граєте раунд 6 з максимум 28.} \\
\texttt{Разом ви накопичили 0 з 16 штрафів.} \\
\texttt{Інший гравець перемістив об'єкт на своїй сітці.} \\
\\ 
\\ 
\texttt{Яка ваша наступна команда?} \\
            }
        }
    }
    & \\ \\

    \theutterance \stepcounter{utterance}  
    & \multicolumn{4}{p{0.6\linewidth}}{
        \cellcolor[rgb]{0.9,0.9,0.9}{
            \makecell[{{p{\linewidth}}}]{
                \texttt{\tiny{[P1$\rangle$GM]}}
                \texttt{СКАЖИ: Добре. Залишилось лише підтвердити готовність: якщо всі об’єкти вже стоять як O(4,1), P(6,1), M(7,1), T(3,2), S(1,4), A(4,5), D(6,7), надішли "завершено?".} \\
            }
        }
    }
    & & \\ \\

    \theutterance \stepcounter{utterance}  
    & & & \multicolumn{2}{p{0.3\linewidth}}{
        \cellcolor[rgb]{0.9,0.9,0.9}{
            \makecell[{{p{\linewidth}}}]{
                \texttt{\tiny{[GM$|$GM]}}
                \texttt{СКАЖИ: Добре. Залишилось лише підтвердити готовність: якщо всі об’єкти вже стоять як O(4,1), P(6,1), M(7,1), T(3,2), S(1,4), A(4,5), D(6,7), надішли "завершено?".} \\
            }
        }
    }
    & & \\ \\

    \theutterance \stepcounter{utterance}  
    & & \multicolumn{4}{p{0.6\linewidth}}{
        \cellcolor[rgb]{0.9,0.9,0.9}{
            \makecell[{{p{\linewidth}}}]{
                \texttt{\tiny{[GM$\rangle$P2]}}
                \texttt{Об'єкт 'S' успішно переміщено в (1,4). Ваша оновлена сітка виглядає так:} \\
\texttt{```} \\
\texttt{ 1234567} \\
\texttt{╔════╤══╗} \\
\texttt{║◌◌◌O│PM║ 1} \\
\texttt{║◌◌T◌├──╢ 2} \\
\texttt{║◌◌◌◌│◌◌║ 3} \\
\texttt{║S┌──┼──╢ 4} \\
\texttt{║A│◌◌│◌◌║ 5} \\
\texttt{╟─┴┬─┤◌D║ 6} \\
\texttt{║◌◌│◌│◌◌║ 7} \\
\texttt{╚══╧═╧══╝} \\
\texttt{```} \\
\texttt{Зараз ви граєте раунд 7 з максимум 28.} \\
\texttt{Разом ви накопичили 0 з 16 штрафів.} \\
\texttt{Інший гравець надіслав це повідомлення:} \\
\\ 
\\ 
\texttt{"Добре. Залишилось лише підтвердити готовність: якщо всі об’єкти вже стоять як O(4,1), P(6,1), M(7,1), T(3,2), S(1,4), A(4,5), D(6,7), надішли "завершено?"."} \\
\\ 
\\ 
\texttt{Яка ваша наступна команда?} \\
            }
        }
    }
    & \\ \\

    \theutterance \stepcounter{utterance}  
    & & & \multicolumn{4}{p{0.6\linewidth}}{
        \cellcolor[rgb]{0.9,0.9,0.9}{
            \makecell[{{p{\linewidth}}}]{
                \texttt{\tiny{[GM$\langle$P2]}}
                \texttt{ПЕРЕСУНЬ: A, (4, 5)} \\
            }
        }
    }
     \\ \\

    \theutterance \stepcounter{utterance}  
    & & & \multicolumn{2}{p{0.3\linewidth}}{
        \cellcolor[rgb]{0.9,0.9,0.9}{
            \makecell[{{p{\linewidth}}}]{
                \texttt{\tiny{[GM$|$GM]}}
                \texttt{ПЕРЕСУНЬ: A, (4, 5)} \\
            }
        }
    }
    & & \\ \\

    \theutterance \stepcounter{utterance}  
    & & & \multicolumn{2}{p{0.3\linewidth}}{
        \cellcolor[rgb]{0.9,0.9,0.9}{
            \makecell[{{p{\linewidth}}}]{
                \texttt{\tiny{[GM$|$GM]}}
                \texttt{A moved to (4, 5)} \\
            }
        }
    }
    & & \\ \\

    \theutterance \stepcounter{utterance}  
    & & \multicolumn{4}{p{0.6\linewidth}}{
        \cellcolor[rgb]{0.9,0.9,0.9}{
            \makecell[{{p{\linewidth}}}]{
                \texttt{\tiny{[P1$\langle$GM]}}
                \texttt{Ваше повідомлення передано іншому гравцю.} \\
\\ 
\\ 
\texttt{Зараз ви граєте раунд 7 з максимум 28.} \\
\texttt{Разом ви накопичили 0 з 16 штрафів.} \\
\texttt{Інший гравець перемістив об'єкт на своїй сітці.} \\
\\ 
\\ 
\texttt{Яка ваша наступна команда?} \\
            }
        }
    }
    & \\ \\

    \theutterance \stepcounter{utterance}  
    & \multicolumn{4}{p{0.6\linewidth}}{
        \cellcolor[rgb]{0.9,0.9,0.9}{
            \makecell[{{p{\linewidth}}}]{
                \texttt{\tiny{[P1$\rangle$GM]}}
                \texttt{СКАЖИ: Чи це вже останній хід? Якщо так і всі 7 об’єктів на цільових місцях, надішли мені "завершено?".} \\
            }
        }
    }
    & & \\ \\

    \theutterance \stepcounter{utterance}  
    & & & \multicolumn{2}{p{0.3\linewidth}}{
        \cellcolor[rgb]{0.9,0.9,0.9}{
            \makecell[{{p{\linewidth}}}]{
                \texttt{\tiny{[GM$|$GM]}}
                \texttt{СКАЖИ: Чи це вже останній хід? Якщо так і всі 7 об’єктів на цільових місцях, надішли мені "завершено?".} \\
            }
        }
    }
    & & \\ \\

    \theutterance \stepcounter{utterance}  
    & & \multicolumn{4}{p{0.6\linewidth}}{
        \cellcolor[rgb]{0.9,0.9,0.9}{
            \makecell[{{p{\linewidth}}}]{
                \texttt{\tiny{[GM$\rangle$P2]}}
                \texttt{Об'єкт 'A' успішно переміщено в (4,5). Ваша оновлена сітка виглядає так:} \\
\texttt{```} \\
\texttt{ 1234567} \\
\texttt{╔════╤══╗} \\
\texttt{║◌◌◌O│PM║ 1} \\
\texttt{║◌◌T◌├──╢ 2} \\
\texttt{║◌◌◌◌│◌◌║ 3} \\
\texttt{║S┌──┼──╢ 4} \\
\texttt{║◌│◌A│◌◌║ 5} \\
\texttt{╟─┴┬─┤◌D║ 6} \\
\texttt{║◌◌│◌│◌◌║ 7} \\
\texttt{╚══╧═╧══╝} \\
\texttt{```} \\
\texttt{Зараз ви граєте раунд 8 з максимум 28.} \\
\texttt{Разом ви накопичили 0 з 16 штрафів.} \\
\texttt{Інший гравець надіслав це повідомлення:} \\
\\ 
\\ 
\texttt{"Чи це вже останній хід? Якщо так і всі 7 об’єктів на цільових місцях, надішли мені "завершено?"."} \\
\\ 
\\ 
\texttt{Яка ваша наступна команда?} \\
            }
        }
    }
    & \\ \\

    \theutterance \stepcounter{utterance}  
    & & & \multicolumn{4}{p{0.6\linewidth}}{
        \cellcolor[rgb]{0.9,0.9,0.9}{
            \makecell[{{p{\linewidth}}}]{
                \texttt{\tiny{[GM$\langle$P2]}}
                \texttt{ПЕРЕСУНЬ: D, (6, 7)} \\
            }
        }
    }
     \\ \\

    \theutterance \stepcounter{utterance}  
    & & & \multicolumn{2}{p{0.3\linewidth}}{
        \cellcolor[rgb]{0.9,0.9,0.9}{
            \makecell[{{p{\linewidth}}}]{
                \texttt{\tiny{[GM$|$GM]}}
                \texttt{ПЕРЕСУНЬ: D, (6, 7)} \\
            }
        }
    }
    & & \\ \\

    \theutterance \stepcounter{utterance}  
    & & & \multicolumn{2}{p{0.3\linewidth}}{
        \cellcolor[rgb]{0.9,0.9,0.9}{
            \makecell[{{p{\linewidth}}}]{
                \texttt{\tiny{[GM$|$GM]}}
                \texttt{D moved to (6, 7)} \\
            }
        }
    }
    & & \\ \\

    \theutterance \stepcounter{utterance}  
    & & \multicolumn{4}{p{0.6\linewidth}}{
        \cellcolor[rgb]{0.9,0.9,0.9}{
            \makecell[{{p{\linewidth}}}]{
                \texttt{\tiny{[P1$\langle$GM]}}
                \texttt{Ваше повідомлення передано іншому гравцю.} \\
\\ 
\\ 
\texttt{Зараз ви граєте раунд 8 з максимум 28.} \\
\texttt{Разом ви накопичили 0 з 16 штрафів.} \\
\texttt{Інший гравець перемістив об'єкт на своїй сітці.} \\
\\ 
\\ 
\texttt{Яка ваша наступна команда?} \\
            }
        }
    }
    & \\ \\

    \theutterance \stepcounter{utterance}  
    & \multicolumn{4}{p{0.6\linewidth}}{
        \cellcolor[rgb]{0.9,0.9,0.9}{
            \makecell[{{p{\linewidth}}}]{
                \texttt{\tiny{[P1$\rangle$GM]}}
                \texttt{СКАЖИ: Схоже, ти вже зробив 7 переміщень. Якщо тепер усе стоїть точно на цілі, надішли "завершено?".} \\
            }
        }
    }
    & & \\ \\

    \theutterance \stepcounter{utterance}  
    & & & \multicolumn{2}{p{0.3\linewidth}}{
        \cellcolor[rgb]{0.9,0.9,0.9}{
            \makecell[{{p{\linewidth}}}]{
                \texttt{\tiny{[GM$|$GM]}}
                \texttt{СКАЖИ: Схоже, ти вже зробив 7 переміщень. Якщо тепер усе стоїть точно на цілі, надішли "завершено?".} \\
            }
        }
    }
    & & \\ \\

    \theutterance \stepcounter{utterance}  
    & & \multicolumn{4}{p{0.6\linewidth}}{
        \cellcolor[rgb]{0.9,0.9,0.9}{
            \makecell[{{p{\linewidth}}}]{
                \texttt{\tiny{[GM$\rangle$P2]}}
                \texttt{Об'єкт 'D' успішно переміщено в (6,7). Ваша оновлена сітка виглядає так:} \\
\texttt{```} \\
\texttt{ 1234567} \\
\texttt{╔════╤══╗} \\
\texttt{║◌◌◌O│PM║ 1} \\
\texttt{║◌◌T◌├──╢ 2} \\
\texttt{║◌◌◌◌│◌◌║ 3} \\
\texttt{║S┌──┼──╢ 4} \\
\texttt{║◌│◌A│◌◌║ 5} \\
\texttt{╟─┴┬─┤◌◌║ 6} \\
\texttt{║◌◌│◌│D◌║ 7} \\
\texttt{╚══╧═╧══╝} \\
\texttt{```} \\
\texttt{Зараз ви граєте раунд 9 з максимум 28.} \\
\texttt{Разом ви накопичили 0 з 16 штрафів.} \\
\texttt{Інший гравець надіслав це повідомлення:} \\
\\ 
\\ 
\texttt{"Схоже, ти вже зробив 7 переміщень. Якщо тепер усе стоїть точно на цілі, надішли "завершено?"."} \\
\\ 
\\ 
\texttt{Яка ваша наступна команда?} \\
            }
        }
    }
    & \\ \\

    \theutterance \stepcounter{utterance}  
    & & & \multicolumn{4}{p{0.6\linewidth}}{
        \cellcolor[rgb]{0.9,0.9,0.9}{
            \makecell[{{p{\linewidth}}}]{
                \texttt{\tiny{[GM$\langle$P2]}}
                \texttt{СКАЖИ: завершено?} \\
            }
        }
    }
     \\ \\

    \theutterance \stepcounter{utterance}  
    & & & \multicolumn{2}{p{0.3\linewidth}}{
        \cellcolor[rgb]{0.9,0.9,0.9}{
            \makecell[{{p{\linewidth}}}]{
                \texttt{\tiny{[GM$|$GM]}}
                \texttt{СКАЖИ: завершено?} \\
            }
        }
    }
    & & \\ \\

    \theutterance \stepcounter{utterance}  
    & & \multicolumn{4}{p{0.6\linewidth}}{
        \cellcolor[rgb]{0.9,0.9,0.9}{
            \makecell[{{p{\linewidth}}}]{
                \texttt{\tiny{[P1$\langle$GM]}}
                \texttt{Ваше повідомлення передано іншому гравцю.} \\
\\ 
\\ 
\texttt{Зараз ви граєте раунд 9 з максимум 28.} \\
\texttt{Разом ви накопичили 0 з 16 штрафів.} \\
\texttt{Інший гравець надіслав це повідомлення:} \\
\\ 
\\ 
\texttt{"завершено?"} \\
\\ 
\\ 
\texttt{Яка ваша наступна команда?} \\
            }
        }
    }
    & \\ \\

    \theutterance \stepcounter{utterance}  
    & \multicolumn{4}{p{0.6\linewidth}}{
        \cellcolor[rgb]{0.9,0.9,0.9}{
            \makecell[{{p{\linewidth}}}]{
                \texttt{\tiny{[P1$\rangle$GM]}}
                \texttt{СКАЖИ: завершено!} \\
            }
        }
    }
    & & \\ \\

    \theutterance \stepcounter{utterance}  
    & & & \multicolumn{2}{p{0.3\linewidth}}{
        \cellcolor[rgb]{0.9,0.9,0.9}{
            \makecell[{{p{\linewidth}}}]{
                \texttt{\tiny{[GM$|$GM]}}
                \texttt{СКАЖИ: завершено!} \\
            }
        }
    }
    & & \\ \\

    \theutterance \stepcounter{utterance}  
    & & & \multicolumn{2}{p{0.3\linewidth}}{
        \cellcolor[rgb]{0.9,0.9,0.9}{
            \makecell[{{p{\linewidth}}}]{
                \texttt{\tiny{[GM$|$GM]}}
                \texttt{true} \\
            }
        }
    }
    & & \\ \\

    \theutterance \stepcounter{utterance}  
    & & & \multicolumn{2}{p{0.3\linewidth}}{
        \cellcolor[rgb]{0.9,0.9,0.9}{
            \makecell[{{p{\linewidth}}}]{
                \texttt{\tiny{[GM$|$GM]}}
                \texttt{Initial states:} \\
\texttt{```} \\
\texttt{Player 1:} \\
\texttt{ 1234567} \\
\texttt{╔════╤══╗} \\
\texttt{║◌◌◌O│PM║ 1} \\
\texttt{║◌◌T◌├──╢ 2} \\
\texttt{║◌◌◌◌│◌◌║ 3} \\
\texttt{║S┌──┼──╢ 4} \\
\texttt{║◌│◌A│◌◌║ 5} \\
\texttt{╟─┴┬─┤◌◌║ 6} \\
\texttt{║◌◌│◌│D◌║ 7} \\
\texttt{╚══╧═╧══╝} \\
\\ 
\\ 
\texttt{Player 2:} \\
\texttt{ 1234567} \\
\texttt{╔════╤══╗} \\
\texttt{║◌◌◌◌│O◌║ 1} \\
\texttt{║P◌M◌├──╢ 2} \\
\texttt{║◌S◌◌│T◌║ 3} \\
\texttt{║◌┌──┼──╢ 4} \\
\texttt{║A│◌◌│◌◌║ 5} \\
\texttt{╟─┴┬─┤◌D║ 6} \\
\texttt{║◌◌│◌│◌◌║ 7} \\
\texttt{╚══╧═╧══╝} \\
\\ 
\texttt{```} \\
            }
        }
    }
    & & \\ \\

    \theutterance \stepcounter{utterance}  
    & & & \multicolumn{2}{p{0.3\linewidth}}{
        \cellcolor[rgb]{0.9,0.9,0.9}{
            \makecell[{{p{\linewidth}}}]{
                \texttt{\tiny{[GM$|$GM]}}
                \texttt{End states:} \\
\texttt{```} \\
\texttt{Player 1:} \\
\texttt{ 1234567} \\
\texttt{╔════╤══╗} \\
\texttt{║◌◌◌O│PM║ 1} \\
\texttt{║◌◌T◌├──╢ 2} \\
\texttt{║◌◌◌◌│◌◌║ 3} \\
\texttt{║S┌──┼──╢ 4} \\
\texttt{║◌│◌A│◌◌║ 5} \\
\texttt{╟─┴┬─┤◌◌║ 6} \\
\texttt{║◌◌│◌│D◌║ 7} \\
\texttt{╚══╧═╧══╝} \\
\\ 
\\ 
\texttt{Player 2:} \\
\texttt{ 1234567} \\
\texttt{╔════╤══╗} \\
\texttt{║◌◌◌O│PM║ 1} \\
\texttt{║◌◌T◌├──╢ 2} \\
\texttt{║◌◌◌◌│◌◌║ 3} \\
\texttt{║S┌──┼──╢ 4} \\
\texttt{║◌│◌A│◌◌║ 5} \\
\texttt{╟─┴┬─┤◌◌║ 6} \\
\texttt{║◌◌│◌│D◌║ 7} \\
\texttt{╚══╧═╧══╝} \\
\\ 
\texttt{```} \\
            }
        }
    }
    & & \\ \\

    \theutterance \stepcounter{utterance}  
    & & & \multicolumn{2}{p{0.3\linewidth}}{
        \cellcolor[rgb]{0.9,0.9,0.9}{
            \makecell[{{p{\linewidth}}}]{
                \texttt{\tiny{[GM$|$GM]}}
                \texttt{* success: True} \\
\texttt{* lose: False} \\
\texttt{* aborted: False} \\
\texttt{{-}{-}{-}{-}{-}{-}{-}} \\
\texttt{* Shifts: 6.00} \\
\texttt{* Max Shifts: 12.00} \\
\texttt{* Min Shifts: 6.00} \\
\texttt{* End Distance Sum: 0.00} \\
\texttt{* Init Distance Sum: 20.21} \\
\texttt{* Expected Distance Sum: 29.33} \\
\texttt{* Penalties: 0.00} \\
\texttt{* Max Penalties: 16.00} \\
\texttt{* Rounds: 9.00} \\
\texttt{* Max Rounds: 28.00} \\
\texttt{* Object Count: 7.00} \\
            }
        }
    }
    & & \\ \\

\end{supertabular}
}

\end{document}
